\section{Design}
\label{sec:design}

\begin{figure}[t]
\centering
\resizebox{\linewidth}{!}{%
\begin{tikzpicture}[
  font=\small,
  box/.style={draw, rounded corners=2pt, align=center, text width=2.6cm, minimum height=3.1em, inner sep=3pt},
  pol/.style={box, fill=red!7}, ws/.style={box, fill=orange!10}, sub/.style={box, fill=blue!6},
  aud/.style={box, fill=green!8},
  arr/.style={-{Stealth[length=2.2mm]}, thick}]
  \node[box] (loop) at (0,0) {agent loop\\\scriptsize unmodified};
  \node[pol] (gate) at (3.7,1.1) {gate\\\scriptsize default-deny, egress ceilings};
  \node[pol] (router) at (3.7,-1.1) {router\\\scriptsize classify outbound payload};
  \node[ws] (ws) at (7.4,0) {working set\\\scriptsize class only rises};
  \node[pol] (engine) at (11.1,-1.1) {placement engine\\\scriptsize rule, ceiling, health};
  \node[sub] (subs) at (11.1,1.1) {substrates\\\scriptsize local · cluster · frontier · network:tool};
  \node[aud] (ledger) at (7.4,-2.6) {ledger\\\scriptsize policy fingerprint, decisions};
  \node[aud] (comply) at (11.1,-2.6) {\texttt{annona comply}\\\scriptsize ledger + policy only};
  \draw[arr] (loop) -- (gate); \draw[arr] (loop) -- (router);
  \draw[arr] (gate) -- (ws); \draw[arr] (router) -- (ws);
  \draw[arr] (ws) -- (engine); \draw[arr] (engine) -- (subs);
  \draw[arr] (gate.east) to[bend left=12] (subs.west);
  \draw[arr] (engine) -- (ledger); \draw[arr] (gate.west) -- ++(-0.3,0) |- (ledger.west);
  \draw[arr] (ledger) -- (comply);
\end{tikzpicture}}
\caption{\annona. Every tool call passes the gate and every inference the router; both fold
what they see into the working set, whose class decides placement. Every decision goes to the
ledger, which a third party audits against the policy alone.}
\label{fig:arch}
\end{figure}

\subsection{Policy}

A policy is a YAML file that must declare \texttt{default: deny}. It defines \emph{classes},
which assign a class to material by path globs and content patterns; \emph{substrates} with
their ceilings; and one \emph{rule} per class. A rule lists the allowed substrates, a
preference among them (privacy, cost, quality or latency, with the listed order as tie-break),
and the action to take when none is available (\texttt{hold}, \texttt{queue}, or a local
\texttt{brief}). Material that matches no class receives the policy's default class, which is
the most restrictive one unless the operator sets it otherwise. Ceilings are declared per
substrate rather than derived from jurisdiction, so an operator can cap an EU substrate at
\pub.

\subsection{The working set}

A tool that runs locally still returns its result into the transcript, and the transcript is
sent with the next inference. For this reason \annona{} keeps a \emph{working set} for each
run. Its class is the maximum over everything the run has touched: paths named by permitted
tool calls, the results of those calls, and paths and patterns found in outgoing payloads. The
working set also keeps provenance, so that \texttt{annona why} can report which file made a run
restricted. The class never decreases. A refused tool call does not raise it; otherwise a
hostile plan could push any run to \res{} by requesting files it will never be allowed to read.

\subsection{Placement}

Before each inference, the router classifies the outgoing payload, adds the result to the
working set, and asks the engine to place a step at the working set's class. The engine selects
the rule for that class and keeps the allowed substrates that can hold the class, support what
the step needs (tools, vision, context length) and are currently available. Each rejected
substrate is recorded with the reason. The remaining candidates are ranked by the rule's
preference. If no candidate remains, the engine applies \texttt{on\_unavailable}; none of its
options adds substrates outside the rule's list. A held step raises an error, the loop returns
the partial result, and the hold is recorded in the ledger like any other outcome.

\subsection{Tools that reach the network}

A gate checks each tool call before it runs. Tools not named in the policy are refused, and a
named tool may only access paths on its allow-list. Deny-lists are evaluated first and
symbolic links are resolved. Our original gate had a gap common to perimeters that only check
inference: the router never saw the arguments of a browser or shell call, so a restricted run
could pass a client file to \texttt{browser} inside a URL. A tool that reaches the network now
has a ceiling. For the shipped browser and shell tools the ceiling is \pub{} unless the policy
raises it. The ceiling is compared with the class of the run, not of the individual call,
because the model may have encoded the material. The ledger records such calls on a
\texttt{network:<tool>} substrate. This label depends only on whether the tool reaches the
network, not on the ceiling applied to it; \cref{sec:eval-compliance} shows why the two must
be kept separate.

\subsection{Content classification}

Path and pattern rules capture where material is stored and which identifiers it contains. They
do not recognise an offer without keywords saved in \texttt{Downloads/}. The classifier
therefore accepts an optional \emph{content model}. It is called wherever content is
classified (a single function used by the router, the gate and result tracking), it can only
raise the class, and an error from it is treated as \res. The included adapter uses the
Jev-compatible \texttt{/v1/systemone} protocol of rizzo-flow~\citep{rizzoflow}. It requests the
probability of each class and compares the \emph{strict mass} $p(\res) + p(\text{cannot tell})$
with a threshold $\tau_\alpha$, set to the $\lfloor \alpha(n+1) \rfloor$-th smallest strict
mass among $n$ labelled restricted texts.

\subsection{Briefs and redaction}

When no available substrate can hold a step's class, the policy may allow a local model to
write a short summary (a \emph{brief}) that is classified again before it can be sent, or may
allow a local redactor (rizzo-pii~\citep{rizzopii}) to replace identifiers while keeping the
mapping on the machine. Material listed under \texttt{sealed} is never summarised or redacted,
because for some documents the sensitive part is the content itself rather than the
identifiers. Both mechanisms are forms of declassification. They are disabled by default for
\res, and we do not claim that they are sound (\cref{sec:discussion}).

\subsection{Ledger and audit}

Each decision is appended to a hash-chained JSON Lines ledger, which stores a digest of the
payload instead of the payload. Each run starts with the SHA-256 fingerprint of the parsed
policy; comments and key order do not affect it, while any change to a rule does. Placement
entries include the class, the rule, the chosen substrate, the remaining candidates and every
rejection. Tool-call entries include the class of the run. \texttt{annona verify} checks the
chain, and \texttt{annona comply} checks the following invariants against the policy file:

\begin{description}[leftmargin=2em, font=\normalfont\ttfamily]
  \item[bound] every run records the fingerprint of this policy;
  \item[sound] every permitted placement uses a substrate that the rule allows and that can hold the class;
  \item[optimal] the chosen substrate is the one ranked first by the rule's preference among the recorded candidates;
  \item[hold-justified] every hold has no recorded candidate;
  \item[monotone] no inference is placed below a class the run has already reached;
  \item[egress] every network tool call happened within its ceiling. The auditor determines
    which tools reach the network from the policy, not from the label written by the daemon.
\end{description}

Checkpoints of the ledger head, published to a witness that the writer cannot modify, detect
truncation and rewritten suffixes before the last checkpoint, as in Certificate Transparency.
